\documentclass{thejus_cv}
\usepackage[hidelinks]{hyperref}

\begin{document}

% --- HEADER ---
\cvheader{Thejus Mahajan Ph.D.}
         {Bioinformatics \& Data Scientist}
         {+49-15259697629}
         {thejus.mahajan@uni-hamburg.de}
         {\href{https://thejusmahajan.github.io}{thejusmahajan.github.io}}
         {Hamburg, Germany}

% Extra contact row for LinkedIn and Google Scholar
\vspace{-0.5cm}
{\small\color{white}
\faLinkedin\ \ \href{https://www.linkedin.com/in/thejusmahajan/}{linkedin.com/in/thejusmahajan}
\hspace{0.5cm}
\faGraduationCap\ \ \href{https://scholar.google.com/citations?user=PJkZwAwAAAAJ}{Google Scholar}
}\par
\vspace{0.1cm}
\begin{paracol}{2}

% ==============================================
% LEFT COLUMN
% ==============================================

\cvsection{Career Interests}
Computational Scientist (PhD) with 5+ years of experience in \textbf{large-scale data analysis}, \textbf{statistical modeling}, and \textbf{pipeline development}. Trained in \textbf{R (Bioconductor, Tidymodels)} and \textbf{Python (Pandas, NumPy, Scikit-learn)} for high-dimensional data analysis. Built a viral genomics pipeline using \textbf{Nextflow (DSL2)}. Currently deepening skills in data pipeline development, \textbf{genomics}, AI/machine learning, and NGS analysis through specialized training and an industry internship. Eager to develop expertise in \textbf{single-cell multi-omics} and apply computational and data science skills to advance research in \textbf{stem cell biology and genomics}.

\cvsection{Experience}

\cvexperience{Internship -- Bioinformatics \& Data Pipeline Engineering}
             {HealthTwiSt GmbH}
             {02/2026 - 04/2026}
             {Berlin, Germany}
             {Clinical data research and development under Dr.~Andreas Busjahn
             \begin{itemize}
                   \item Refactored a monolithic 1,834-line R data pipeline (tidyverse) for Germany's interventional radiology quality registry (DeGIR); 143,000+ patient records from \textasciitilde300 clinics. Reduced to 1,348 lines ($-$26.5\%) with byte-identical output verified via \texttt{identical()}.
                   \item Externalised 257 hard-coded correction rules into 8 configuration CSV files, enabling non-programmer editing of business rules.
                   \item Replaced 357 deprecated function calls across 4 pipeline files, modernising the codebase to current tidyverse standards.
                   \item Created the \texttt{degirtools} R package (9 exported functions, roxygen2 docs, \texttt{devtools::check()} passing) by extracting reusable logic from a 767-line helper file.
                   \item Discovered 2 pre-existing bugs through systematic code analysis; documented both and preserved original behaviour for supervisor review --- demonstrating production-code discipline.
                   \item Built an interactive R Shiny dashboard (6 modules: overview, interventions, complications, radiation doses, success rates, about) using plotly, DT, and bslib. Deployed to shinyapps.io with GDPR-safe synthetic data.
                   \item Wrote Tidymodels teaching materials (Quarto/GitHub) replacing the legacy caret framework for ML workflows in R.
               \end{itemize}}

\cvexperience{Further Training -- Bioinformatics and Biostatistics}
             {CQ Beratung + Bildung GmbH}
             {08/2025 - 02/2026}
             {Berlin, Germany}
             {Application-related Bioinformatics and Biostatistics
             \begin{itemize}\setlength{\itemsep}{0.05em}
                 \item Programming methodology: Shell (BASH), Python (object-oriented), R.
                 \item Bioinformatics databases (NCBI, Biopython, SQL) and sequence analysis.
                 \item NGS analysis pipelines (Nextflow, Galaxy); structural bioinformatics.
                 \item Clinical biostatistics (R/Bioconductor, ANOVA, PCA, HCA, survival analysis); GDPR/BDSG compliance.
                 \item Capstone: Built a reproducible Nextflow (DSL2) pipeline for Hepatitis Delta Virus sequence analysis --- automated NCBI download, MAFFT alignment, and trimAl trimming with Conda-managed dependencies (\href{https://github.com/thejusmahajan/hepatitis-delta-pipeline}{\faGithub\ GitHub}).
             \end{itemize}}

\cvexperience{Guest Scientist}
             {Helmholtz-Zentrum Hereon, \"Okosystemmodellierung}
             {05/2025 - 10/2025}
             {Geesthacht, Germany}
             {Computational Biologist \& Systems Modeler
             \begin{itemize}
                 \item Spearheaded research modeling evolutionary processes in marine ecosystems.
                 \item Applied advanced computational techniques to simulate long-term environmental impacts.
                 \item Finalised the computational modelling framework and prepared the simulation data for peer-reviewed publication.
             \end{itemize}}

\cvexperience{Job Search \& German Language Learning}
             {Hamburg, Germany}
             {02/2025 - 04/2025}
             {Hamburg, Germany}
             {Active job search and German language acquisition (Goethe B1 certification).}

\cvexperience{Post-doctoral Scientist}
             {University of Hamburg, Institute of Marine Ecosystem and Fisheries Sciences}
             {08/2021 - 01/2025}
             {Hamburg, Germany}
             {Marine Ecosystem Modelling
             \begin{itemize}
                 \item Developed the novel ``Cyanobacteria Life Cycle'' (CLC) model within the ERGOM framework for climate warming projections.
                 \item Engineered models in Fortran and analyzed large NetCDF datasets using Python in Linux.
                 \item Statistical analysis and visualization of ecological time-series data using R.
             \end{itemize}}

\cvexperience{Job Search}
             {Kerala, India}
             {04/2021 - 07/2021}
             {Kerala, India}
             {Relocation from India to Germany and active job search.}

\cvexperience{Physics Educator}
             {Tutorwaves Solutions Inc / Freelance}
             {10/2018 - 03/2021}
             {Kerala, India}
             {Prepared physics content for competitive examinations and coached state-level chess teams.}

\cvexperience{Ph.D. Researcher in Astrochemistry}
             {University of Paris-Saclay}
             {10/2015 - 09/2018}
             {Orsay, France}
             {Led atomic collision experiments, analyzed data, and modeled outcomes.
             \begin{itemize}
                 \item Acquired and processed terabytes of raw data from tandem particle accelerator experiments.
                 \item C++ signal processing, algorithm development, and molecular fragmentation modeling (KIDA database).
             \end{itemize}}


\cvsection{Education}

\cvexperience{Ph.D. in Astrochemistry}
             {Universit\'e Paris-Saclay, Institute of Molecular Sciences (ISMO) Orsay, France}
             {09/2015 - 09/2018}
             {Orsay, France}
             {}

\cvexperience{M.Sc. in Physics}
             {National Institute of Technology Calicut, India}
             {07/2012 - 01/2015}
             {Calicut, India}
             {}

\cvexperience{B.Sc. in Physics}
             {University of Calicut, India}
             {06/2009 - 04/2012}
             {Thrissur, India}
             {}

\cvsection{References}

\begin{tabular}{@{}p{0.5\linewidth} p{0.5\linewidth}@{}}
\textbf{Dr. Andreas Busjahn} & \textbf{Prof. Dr. Kai Wirtz} \\
abusjahn@healthtwist.de & kai.wirtz@hereon.de \\[0.4cm]
\textbf{Prof. Dr. Elisa Schaum} & \\
elisa.schaum@uni-hamburg.de & \\
\end{tabular}
\vspace{0.4cm}

\switchcolumn

% ==============================================
% RIGHT COLUMN
% ==============================================

\cvsection{Strengths}

\cvstrength{\faStar}
           {Data Pipeline Engineering}
           {Refactored a production data pipeline (1,834 lines, 143K records, \textasciitilde300 clinics) with 26.5\% code reduction. Built degirtools R package. Byte-level verification after every change via \texttt{identical()}.}

\cvstrength{\faLightbulb}
           {Complex Problem-Solving}
           {Resolved a critical data-loss error by creating a Fortran simulation to correct for a misplaced particle detector, recalculating results using multi-sensor data.}

\cvstrength{\faAward}
           {Bioinformatics \& Biostatistics}
           {Bridged physical science expertise with modern biological workflows. Skilled in NGS data analysis, sequence alignment, and multivariate statistical modeling using R (Bioconductor) and Python.}

\cvsection{Technical Skills}

\cvskills{Programming}
         {Python (Pandas, NumPy, Scikit-learn, Biopython), R (Bioconductor, Tidymodels, ggplot2), SQL, BASH}

\cvskills{Bioinformatics}
         {Nextflow (DSL2), NGS Analysis, Multiple Sequence Alignment (MAFFT), Sequence Trimming (trimAl), NCBI / Entrez}

\cvskills{Biostatistics \& ML}
         {Multivariate Analysis, PCA, ANOVA, Dimensionality Reduction, Clustering, Scikit-learn, Tidymodels}

\cvskills{Data Engineering}
         {Large dataset processing (HDF5, NetCDF), Pipeline development, Data visualization, R Shiny, devtools/roxygen2}

\cvskills{Tools}
         {Linux/Unix, Git, HPC clusters, Conda}


\cvsection{Training / Courses}

\textbf{IBM: Machine Learning with Python; Introduction to AI}\\
{\small Coursera (in progress)}\\[0.3cm]

\textbf{High Performance Computing (HPC) Training}\\
J\"ulich Supercomputing \textbf{Center (JSC)},\\
Forschungszentrum J\"ulich\\[0.3cm]

\textbf{MITx: 7.00x Introduction to Biology - The Secret of Life}\\
{\small Comprehensive coursework on genetics, biochemistry, and molecular biology.}\\[0.4cm]

\cvsection{Other Selected Projects}

\textbf{Hepatitis Delta Virus Pipeline} \textbullet\ \href{https://github.com/thejusmahajan/hepatitis-delta-pipeline}{\faGithub\ GitHub}\par
{\small Nextflow (DSL2) pipeline for viral sequence analysis: NCBI Entrez download, MAFFT alignment, trimAl trimming, Conda-managed dependencies.}\par
\vspace{0.3em}

\textbf{Introduction to Tidymodels} \textbullet\ \href{https://github.com/thejusmahajan/Introduction_to_Tidymodels}{\faGithub\ GitHub}\par
{\small Teaching script for ML in R: recipes, parsnip, workflows, and yardstick on the Palmer Penguins dataset.}\par
\vspace{0.3em}

\textbf{Berlin Tree Analysis} \textbullet\ \href{https://github.com/thejusmahajan/berlin-tree-analysis}{\faGithub\ GitHub}\par
{\small Interactive Jupyter notebook analyzing 5 years of tree planting data in Berlin using Python, Pandas, and Seaborn.}\par
\vspace{0.4em}

\cvsection{Publications}

5 peer-reviewed articles \textbullet\ \href{https://scholar.google.com/citations?user=PJkZwAwAAAAJ}{\faGraduationCap\ Google Scholar Profile}

\cvsection{Languages}

\cvlanguage{English}{Proficient}{4}
\cvlanguage{German}{B1 - Goethe Certified}{3}
\cvlanguage{French}{Intermediate}{2}
\cvlanguage{Malayalam}{Native}{5}
\cvlanguage{Tamil}{Proficient}{4}
\cvlanguage{Hindi}{Intermediate}{2}

\end{paracol}
\end{document}
